\section{Related Work}
\label{app:related-work}

\paragraph{Symbolic regression benchmarks.}
Existing SR benchmarks primarily focus on constructing static task collections and reporting final-performance comparisons under a fixed method set. SRBench established one of the first large-scale reproducible evaluation efforts for contemporary symbolic regression, comparing 14 SR methods and 7 machine-learning baselines across 252 regression problems \citep{lacava2021contemporary}. SRSD revisited formula-backed scientific discovery by reconstructing Feynman-style datasets, adding dummy-variable variants, and introducing normalized edit distance for structural expression comparison \citep{matsubara2024rethinking}. LLM-SRBench targets the emerging LLM-based equation-discovery setting with 239 tasks designed to reduce memorization effects and evaluate scientific reasoning beyond common formulas \citep{shojaee2025llmsrbench}. These resources provide important task pools and evaluation conventions. However, they are still largely static benchmark suites: they do not specify how to distill a large heterogeneous reservoir into a compact core benchmark, how to keep scores stable as new algorithms are added, or how to audit search dynamics through intermediate expressions. \arena{} builds on these benchmark resources, but shifts the focus from dataset collection to evaluation infrastructure.

\paragraph{Limitations of static SR evaluation.}
Most existing SR evaluations report final numerical metrics such as MSE, NMSE, or $R^2$ after a fixed search budget \citep{lacava2021contemporary,matsubara2024rethinking,shojaee2025llmsrbench}. This protocol is useful for one-time comparison, but it misses several properties that are central to modern SR systems. First, low numerical error does not necessarily imply symbolic correctness: a method can fit the test region with a non-equivalent shortcut expression. Second, final metrics hide search dynamics: two algorithms with the same final NMSE may differ substantially in time-to-discovery, early-stage quality, or complexity growth. Third, static benchmark tables do not provide stable scores under leaderboard growth if scores are rank-normalized or computed relative to the current method pool. Adjacent evaluation-infrastructure work has made similar criticisms of static benchmarking in NLP and data-centric AI \citep{kiela2021dynabench,ma2021dynaboard,mazumder2023dataperf}. \arena{} addresses these gaps for SR by logging minute-level best-so-far expressions, canonicalizing outputs into expression trees, and using fixed absolute metrics whose values do not drift when new algorithms are added.

\paragraph{Modern SR algorithms.}
The SR algorithmic landscape has diversified beyond classical genetic programming. Evolutionary and GP-style systems, including PySR, PyOperon, and gplearn, search over expression trees and often combine structure search with constant optimization \citep{cranmer2023pysr}. Physics-inspired systems such as AI Feynman exploit inductive biases including symmetry, separability, and compositional structure to decompose equation discovery into simpler subproblems \citep{udrescu2020aifeynman}. Reinforcement-learning and policy-search approaches, such as DSR, learn expression-generating policies using reward signals derived from numerical fit and expression quality \citep{petersen2021dsr}. Neural-guided and hybrid methods combine learned proposal mechanisms with symbolic search or genetic-programming components \citep{mundhenk2021dsoseeding}. Transformer-based and planning-based systems use pretraining or tree search to guide symbolic generation \citep{shojaee2023tpsr}. More recent LLM-assisted SR systems use language models to propose equation programs, scientific hypotheses, or context-aware symbolic structures \citep{shojaee2025llmsr}. Because these paradigms differ in prior usage, search trajectory, symbolic faithfulness, runtime cost, OOD behavior, and seed stability, a single final-error table is insufficient to characterize them. \arena{} is not a new SR algorithm; it is designed to evaluate this increasingly heterogeneous algorithmic space under one executable contract.

\paragraph{LLM-assisted equation discovery.}
LLM-assisted SR introduces a new evaluation challenge. Unlike traditional symbolic search, LLM-based methods can exploit natural-language descriptions, scientific priors, memorized formula patterns, or prompt-specific reasoning strategies \citep{shojaee2025llmsr,shojaee2025llmsrbench}. This makes direct comparison with non-LLM systems difficult unless the evaluation protocol explicitly controls what contextual information is provided. LLM-oriented benchmarks therefore emphasize anti-memorization design and scientific equation-discovery tasks \citep{shojaee2025llmsrbench}. \arena{} separates these concerns in two ways. First, PySR and LLM-SR are used as dual probes without physical background information during \candidate{} mining, so the initial high-information pool is not merely a test of language-context advantage. Second, LLM-assisted methods can still participate in the formal leaderboard, but their context usage and prompt contract are explicitly recorded. This separation allows \arena{} to use LLM systems as informative probes without letting LLM priors dominate the construction of the final \core{}.

\paragraph{Benchmark distillation and core-set selection.}
Large benchmark suites improve coverage, but their cost makes them difficult to use during iterative algorithm development. A common workaround is to evaluate on a smaller subset, but naive random sampling, family-stratified sampling, or difficulty-only selection can miss the behavioral conclusions induced by the full method panel. Recent work on efficient model evaluation and benchmark subset selection argues that smaller benchmarks should preserve more than item diversity: they should also preserve method ranking, item discrimination, and distributional coverage \citep{perlitz2024efficient,polo2024tinybenchmarks,vivek2024anchorpoints,saranathan2025sublime}. \arena{} treats subset construction as benchmark distillation rather than convenience sampling. Its \core{} objective jointly considers structural coverage, response-space coverage, method discriminability, redundancy control, difficulty and failure-mode balance, and seed-level stability. This differs from selecting the easiest, hardest, or largest-gap tasks. It also differs from metadata-only core-set selection because \arena{} uses probe-algorithm responses to preserve behavioral conclusions from the larger 664-task reservoir.

\paragraph{Probe-based evaluation and algorithm-panel fidelity.}
A key challenge in benchmark distillation is that the selected subset should represent not only the task distribution, but also the conclusions that a larger algorithm panel would induce. \arena{} addresses this through a cascaded probe design. First, PySR and LLM-SR are used as dual probes to mine a high-information \candidate{} pool, intentionally exposing differences between a strong traditional SR system and an LLM-assisted paradigm. Second, \arena{} evaluates 12 integrated algorithms on \candidate{} and selects four non-LLM probes through a health-gated exhaustive subset-selection procedure. The selected \probe{} is optimized for health, panel fidelity, behavioral complementarity, and dataset coverage rather than average performance alone. This probe-selection step is important because the final \core{} is not chosen to favor the strongest algorithms; it is chosen to approximate the dataset-informativeness judgments of a broader algorithm panel while remaining computationally tractable.

\paragraph{Expression-level evaluation.}
Several SR benchmarks recognize that numerical error alone is insufficient for scientific discovery. SRSD, for example, introduced normalized edit distance to compare predicted expressions with ground-truth formulas \citep{matsubara2024rethinking}. This is important because symbolic regression is often used not merely to predict values, but to recover interpretable mechanisms. \arena{} extends expression-level evaluation in two directions. First, all outputs are canonicalized into a common expression-tree representation, allowing equivalence testing, tree similarity, variable recovery, operator recovery, and seed-level structural consistency. Second, symbolic fidelity is integrated into the leaderboard as a first-class axis rather than treated as a post-hoc diagnostic. This makes it possible to distinguish low-NMSE equivalent formulas from low-NMSE but symbolically incorrect shortcuts.

\paragraph{OOD generalization and robustness in SR.}
OOD behavior is central to equation discovery because a formula that only interpolates the training domain may not reflect the underlying mechanism. Prior SR benchmarks have considered extrapolation or scientific-discovery splits to varying degrees \citep{lacava2021contemporary,matsubara2024rethinking,shojaee2025llmsrbench}. However, OOD evaluation is often reported as another final-error column, without explicitly separating absolute OOD quality from ID-to-OOD degradation. \arena{} defines \oodg{} as a combination of OOD quality and ID-to-OOD retention, preventing methods from being rewarded simply because they perform poorly both in and out of distribution. \arena{} further defines \rob{} as a noisy-training extension track, measuring whether algorithms trained on corrupted labels can still recover clean ID/OOD behavior. This separates clean-track evaluation from robustness evaluation while keeping both under a shared metric protocol.

\paragraph{Efficiency, anytime behavior, and search traces.}
SR methods can differ substantially in how quickly they discover useful expressions. Evolutionary systems may gradually improve over time, neural-guided systems may find strong candidates early, and LLM-assisted systems may generate high-quality initial structures but vary in refinement behavior. Existing SR benchmarks typically report final scores after a fixed budget, which hides these anytime differences. Dynamic benchmarking and evaluation-as-a-service work in other domains has argued for richer evaluation artifacts, stable infrastructure, and additional metrics beyond accuracy \citep{kiela2021dynabench,ma2021dynaboard}. \arena{} brings this idea to SR by recording best-so-far expressions every minute under a shared one-hour budget, enabling efficiency metrics based on quality-time AUC, time-to-threshold, stagnation analysis, and complexity-over-time. This turns the leaderboard from a final-score table into a diagnostic tool for understanding search behavior.

\paragraph{Evaluation infrastructure and dataset documentation.}
Recent evaluation practice increasingly emphasizes not only benchmark results, but also artifact documentation, dataset cards, model cards, data statements, metadata, and reproducibility pathways \citep{gebru2021datasheets,mitchell2019modelcards,bender2018datastatements,pushkarna2022datacards,akhtar2024croissant}. Data-centric evaluation efforts such as DataPerf also argue that benchmark infrastructures should support iterative development and long-term maintenance rather than one-off static comparisons \citep{mazumder2023dataperf}. \arena{} follows this direction by releasing dataset manifests, standardized result schemas, expression canonicalization tools, metric-computation scripts, license tables, and Croissant-style metadata. Its contribution is therefore not only the \core{} task set, but the infrastructure connecting task standardization, algorithm wrappers, result semantics, benchmark distillation, dynamic leaderboard updates, and expression-level diagnostics. In this sense, \arena{} aims to make SR benchmarking auditable and extensible rather than merely rank-based.
